\chapter{Installation und erster Start}\label{kap:installation}
\index{Installation}

Die Installation hat zwei Teile: das Programm selbst, das mit allen mitgelieferten Plugins
in einem Paket kommt, und eine Python-Umgebung, in der die Plugins rechnen. Das Programm
allein genügt zum Modellieren, Betrachten und Konvertieren von Geometrie; für Berechnungen
brauchen Sie zusätzlich die Python-Umgebung (\autoref{sec:python-umgebung}).

Die Installationspakete stehen auf der Seite der Veröffentlichungen (\enquote{Releases}) des
Projekts:
\begin{lstlisting}
https://github.com/PhotonPower/Raytatouille-Fusion/releases
\end{lstlisting}

\section{Systemvoraussetzungen}\label{sec:systemvoraussetzungen}
\index{Systemvoraussetzungen}

\begin{referenz}{>{\raggedright\arraybackslash}p{3.3cm}>{\raggedright\arraybackslash}p{5.4cm}X}
\toprule
\textbf{System} & \textbf{Paket} & \textbf{Hinweise} \\
\midrule
\endhead
macOS 14 (Sonoma) oder neuer, Apple Silicon und Intel &
  \datei{RaytatouilleFusion-<Version>-macos-universal.dmg} &
  Ein Programm für beide Prozessorarten (\enquote{universal}). Nicht signiert, siehe
  \autoref{sec:inst-macos}. Grafik über Metal. \\
Windows 11 (64 Bit) &
  \datei{RaytatouilleFusion-<Version>-win64-setup.exe} &
  Installationsprogramm; installiert für alle Benutzer des Rechners. Grafik über
  Direct3D~11. \\
AlmaLinux 10 / Red Hat Enterprise Linux 10 (x86\_64) &
  \datei{raytatouille-fusion-<Version>-<Release>.x86\_64.rpm} &
  Nutzt die Qt-6-Pakete der Distribution. Grafik über OpenGL (Vulkan wählbar). Läuft unter
  Wayland und X11. \\
\bottomrule
\end{referenz}

Weitere Voraussetzungen:

\begin{itemize}
  \item \textbf{Grafik:} eine Grafikkarte bzw.\ ein Grafiktreiber mit Unterstützung für die
    genannte Schnittstelle. In virtuellen Maschinen hilft gegebenenfalls OpenGL
    (\autoref{sec:umgebungsvariablen}).
  \item \textbf{Python} für die Plugins: Version~3.10 oder neuer ist Pflicht -- ältere
    Versionen meldet \gui{Check environment} als \enquote{too old}. Empfohlen werden die
    Versionen~3.10 bis~3.13, am besten~3.13; für sie liegen die Rechenkerne als fertig
    übersetzte Pakete (\enquote{Wheels}) für Windows, Linux und macOS bereit, sodass kein
    Compiler nötig ist.
  \item \textbf{Internet} beim Einrichten der Python-Umgebung (die Pakete kommen aus dem
    Python Package Index, PyPI) und für die Update-Prüfung. Das Programm selbst arbeitet
    ohne Netz und sendet von sich aus nichts.
  \item \textbf{Arbeitsspeicher und Prozessor:} Zum Modellieren genügt ein üblicher
    Arbeitsplatzrechner. Rechnungen der Plugins (Kamerabilder mit vielen Strahlen, Path
    Tracing, Wellenoptik) profitieren von vielen Prozessorkernen und viel Arbeitsspeicher;
    der Workflow behält Zwischenergebnisse bis zu \SI{1}{GB}.
\end{itemize}

\section{Installation unter macOS}\label{sec:inst-macos}
\index{Installation!macOS}\index{macOS}

\begin{schritte}
  \item Laden Sie die Datei \datei{RaytatouilleFusion-<Version>-macos-universal.dmg}
    herunter und öffnen Sie sie mit einem Doppelklick. Es erscheint ein Fenster
    \enquote{Raytatouille Fusion} mit dem Programm \datei{RaytatouilleFusion} und einer
    Verknüpfung zum Ordner \datei{Applications} (Programme).
  \item Ziehen Sie \datei{RaytatouilleFusion} auf \datei{Applications}.
  \item Werfen Sie das Disk-Image aus (Papierkorb-Symbol neben dem Laufwerk im Finder).
  \item Öffnen Sie das Programm beim \textbf{ersten} Mal nicht per Doppelklick, sondern mit
    \maus{Rechtsklick} (oder \keys{Ctrl}-Klick) auf das Programm im Ordner Programme
    $\to$ \gui{Öffnen}, und bestätigen Sie die Rückfrage mit \gui{Öffnen}.
\end{schritte}

\begin{hinweis}
Das Programm ist bewusst nicht von Apple signiert und notarisiert. macOS warnt deshalb beim
ersten Start, dass der Entwickler nicht überprüft werden kann. Bietet der Dialog keinen Knopf
zum Öffnen an, gehen Sie in die \gui{Systemeinstellungen} $\to$ \gui{Datenschutz \&
Sicherheit}; dort steht unten eine Meldung zu \enquote{RaytatouilleFusion} mit dem Knopf
\gui{Dennoch öffnen}. Diese Freigabe ist nur einmal nötig; danach startet das Programm normal
per Doppelklick, aus dem Launchpad oder über Spotlight.
\end{hinweis}

\paragraph{Wo liegt was?} Im Ordner Programme (\datei{/Applications}) liegt das Programm als
ein einziges Programmpaket. Darin liegen auch die mitgelieferten Plugins und das Python-Paket
\code{nfcad}, über das die Plugins mit dem Programm sprechen. Die ausführbare Datei für den
Start aus dem Terminal ist
\begin{lstlisting}
/Applications/RaytatouilleFusion.app/Contents/MacOS/RaytatouilleFusion
\end{lstlisting}

\paragraph{Update.} Ziehen Sie die neue Version ebenso nach \datei{Applications} und lassen
Sie die alte ersetzen. Einstellungen und Plugin-Umgebung bleiben erhalten.

\section{Installation unter Windows 11}\label{sec:inst-windows}
\index{Installation!Windows}\index{Windows}

\begin{schritte}
  \item Laden Sie \datei{RaytatouilleFusion-<Version>-win64-setup.exe} herunter und starten Sie
    die Datei.
  \item Windows fragt nach Administratorrechten, weil das Programm für alle Benutzer nach
    \datei{C:\textbackslash Program Files\textbackslash Raytatouille Fusion} installiert
    wird. Bestätigen Sie.
  \item Erscheint die Meldung \enquote{Der Computer wurde durch Windows geschützt}
    (Microsoft Defender SmartScreen)\index{SmartScreen}, klicken Sie auf \gui{Weitere
    Informationen} und dann auf \gui{Trotzdem ausführen}. Die Warnung erscheint, weil das
    Installationsprogramm nicht signiert ist.
  \item Folgen Sie dem (englischsprachigen) Assistenten: Lizenz (MIT) annehmen und, falls
    gefragt, Zielordner und Startmenü-Ordner bestätigen. Auf der Seite mit den zusätzlichen
    Aufgaben können Sie mit \gui{Create a desktop shortcut} ein Symbol auf dem Desktop
    anlegen lassen.
  \item Auf der letzten Seite startet \gui{Launch Raytatouille Fusion} das Programm direkt.
\end{schritte}

Das Installationsprogramm
\begin{itemize}
  \item legt einen Eintrag \gui{Raytatouille Fusion} im Startmenü an (und auf Wunsch auf dem
    Desktop),
  \item verknüpft die Dateiendungen \datei{.rttf} und \datei{.nfcad} mit dem Programm, sodass
    ein Doppelklick auf ein Dokument es öffnet (\autoref{sec:kommandozeile}),
  \item installiert das Programm als \datei{bin\textbackslash RaytatouilleFusion.exe} und die
    mitgelieferten Plugins unter \datei{python\textbackslash plugins} im Zielordner.
\end{itemize}

\paragraph{Update.} Führen Sie einfach das Installationsprogramm der neuen Version aus. Es
erkennt die vorhandene Installation und ersetzt sie; das gilt auch für eine Installation des
Vorgängers \enquote{Nurbforge CAD}.

\section{Installation unter AlmaLinux und RHEL 10}\label{sec:inst-linux}
\index{Installation!Linux}\index{Linux}

Installieren Sie das RPM-Paket mit \code{dnf}; es holt die benötigten Qt-6-Pakete der
Distribution (\code{qt6-qtbase-gui}, \code{qt6-qtsvg}) automatisch dazu und empfiehlt
\code{python3}:

\begin{lstlisting}
sudo dnf install ./raytatouille-fusion-*.rpm
\end{lstlisting}

Das Paket installiert
\begin{itemize}
  \item das Programm \datei{/usr/bin/RaytatouilleFusion},
  \item einen Eintrag im Anwendungsmenü (\gui{Raytatouille Fusion}, Kategorie Grafik /
    Technik) mit Symbol und AppStream-Metadaten,
  \item die mitgelieferten Plugins unter \datei{/usr/share/raytatouille-fusion/python}.
\end{itemize}

Starten Sie das Programm aus dem Anwendungsmenü oder im Terminal mit dem
Befehl \code{RaytatouilleFusion}. Das Paket hieß früher \code{nurbforge-cad}; es wird beim
Installieren ersetzt.

\begin{tipp}
Das Programm läuft unter Wayland und X11. Gibt es unter Wayland Darstellungsprobleme,
erzwingen Sie X11 mit \code{QT\_QPA\_PLATFORM=xcb RaytatouilleFusion}.
\end{tipp}

\paragraph{Update.} Installieren Sie das neue Paket ebenso:
\begin{lstlisting}
sudo dnf upgrade ./raytatouille-fusion-*.rpm
\end{lstlisting}

\section{Die Python-Umgebung für die Plugins}\label{sec:python-umgebung}
\index{Python-Umgebung}

\subsection{Warum eine eigene Python-Umgebung?}

Die Plugins sind Python-Programme. Sie werden vom CAD-Programm als eigene Prozesse mit
\emph{Ihrem} Python-Interpreter gestartet und brauchen dort die Rechenkerne als
Python-Pakete. Mitgeliefert werden nur die Plugins selbst und das kleine Paket \code{nfcad},
nicht aber Python und die Rechenkerne. Das hat zwei Vorteile: Sie können die Rechenkerne
unabhängig vom Programm aktualisieren, und Sie können dieselbe Umgebung auch in eigenen
Skripten verwenden.

Die Plugins brauchen zusammen die folgenden Pakete (sie stehen in den Dateien
\datei{plugin.json} der Plugins unter \param{requires}):

\begin{referenz}{>{\raggedright\arraybackslash}p{4.8cm}>{\raggedright\arraybackslash}p{4.2cm}X}
\toprule
\textbf{Paket (pip-Name)} & \textbf{Mindestversion} & \textbf{gebraucht von} \\
\midrule
\endhead
\code{numpy} & -- & allen Plugins und \code{nfcad} \\
\code{raytatouille} & 0.9 & Raytracer, Workflow \\
\code{rtt-fusion} (Import \code{rtt\_fusion}) & 0.2 & Raytracer, Workflow \\
\code{photonpower-fieldtracing} (Import \code{fieldtracing}) & 0.2 & Raytracer \\
\code{fouriergrating} & 0.18 & FourierGrating \\
\code{scipy} & -- & FourierGrating \\
\code{hpfem} & 0.5 & hp-FEM \\
\code{pathtracer} & 0.3 & Path Tracer, Automotive Lighting \\
\bottomrule
\end{referenz}

\subsection{Schritt 1: Python installieren}

Installieren Sie einmalig ein Python ab Version~3.10, empfohlen~3.13:

\begin{description}
  \item[Windows:] Installationsprogramm von \url{https://www.python.org/downloads/}. Setzen Sie
    im ersten Bild des Installers das Häkchen bei \gui{Add python.exe to PATH}.
  \item[macOS:] Installationspaket von python.org oder über Homebrew:
\begin{lstlisting}
brew install python@3.13
\end{lstlisting}
  \item[AlmaLinux/RHEL:] \code{sudo dnf install python3.12} (oder eine andere angebotene
    Version ab 3.10).
\end{description}

\begin{hinweis}
\textbf{Welcher Python-Interpreter wird benutzt?} Das Programm nimmt der Reihe nach:
\begin{enumerate}
  \item den unter \menu{Plugins,Settings...} eingetragenen Interpreter,
  \item sonst den aus der Umgebungsvariablen \code{NFCAD\_PYTHON},
  \item sonst das erste \code{python3} oder \code{python} im Suchpfad (\code{PATH}); unter
    Windows \code{python}, \code{py} oder \code{python3}, wobei die Platzhalter des
    Microsoft Store (unter \datei{\textbackslash WindowsApps\textbackslash}) übersprungen
    werden -- sie würden nur den Store öffnen.
\end{enumerate}
Unter macOS sieht ein aus dem Finder gestartetes Programm nicht unbedingt denselben
Suchpfad wie das Terminal; tragen Sie den Interpreter im Zweifel ausdrücklich ein. Am
sichersten ist eine eigene virtuelle Umgebung, wie sie der nächste Schritt anlegt.
\end{hinweis}

\subsection{Schritt 2: Umgebung anlegen mit \enquote{Create environment...}}
\label{sec:create-environment}\index{Create environment}

Das Programm legt auf Knopfdruck eine virtuelle Umgebung (\enquote{venv}) an und installiert
alle Pakete, die die geladenen Plugins brauchen:

\begin{schritte}
  \item Wählen Sie \menu{Plugins,Settings...}. Es öffnet sich der Dialog \gui{Plugin
    settings} (\autoref{fig:plugin-settings}).
  \item Prüfen Sie das Feld \gui{Python interpreter}. Es muss auf das in Schritt~1
    installierte Python zeigen (mit dem Knopf \gui{...} daneben können Sie eine Datei
    auswählen). Dieses Python dient nur als Vorlage für die neue Umgebung.
  \item Klicken Sie auf \gui{Create environment...}. Ein Dateidialog \gui{New Python
    environment (folder)} fragt nach dem Ordner der neuen Umgebung. Vorgeschlagen wird der
    Ordner \datei{python-env} im Datenordner des Programms (\autoref{sec:einstellungen}).
    Der Ordner darf noch \emph{nicht} existieren; sonst erscheint \enquote{\dots{} exists
    already -- choose a new folder.}
  \item Das Programm führt nun zwei Schritte aus und schreibt ihre Ausgabe in das Textfeld
    darunter: zuerst \code{python -m venv <Ordner>} (\enquote{Creating \dots{} from \dots}),
    dann \code{pip install} mit allen benötigten Paketen (\enquote{Installing numpy, \dots}).
    Das kann einige Minuten dauern. Der Dialog bleibt dabei bedienbar.
  \item Am Ende steht \enquote{Done -- the new interpreter is set above (applied with OK).}
    Das Feld \gui{Python interpreter} zeigt jetzt den Interpreter der neuen
    Umgebung: \datei{bin/python} im gewählten Ordner, unter Windows
    \datei{Scripts\textbackslash python.exe}.
  \item Klicken Sie auf \gui{Check environment}. Im Textfeld sollte die Python-Version mit
    \enquote{ok} und vor jedem Paket ein Häkchen (\checkmark) stehen.
  \item Bestätigen Sie den Dialog mit \gui{OK}. Erst dann wird der neue Interpreter
    übernommen; die Plugins werden neu gestartet.
\end{schritte}

\abbildung[0.85\linewidth]{plugin-settings.png}{Der Dialog \gui{Plugin settings}
(\menu{Plugins,Settings...}) mit dem Interpreter, den Knöpfen für die Umgebung, dem Bericht
von \gui{Check environment}, den Plugin-Ordnern und dem Entwicklermodus. Die Aufnahme stammt
aus einer Entwicklungsumgebung; Pfade und Pakete sehen bei Ihnen anders aus.}{fig:plugin-settings}

\begin{warnung}
Wenn Sie den Dialog mit \gui{Cancel} schließen, wird der neue Interpreter \emph{nicht}
übernommen, obwohl die Umgebung angelegt ist. Öffnen Sie den Dialog dann erneut, tragen Sie
den Interpreter der Umgebung ein (\gui{...}) und bestätigen Sie mit \gui{OK}.
\end{warnung}

Schlägt die Installation der Pakete fehl (\enquote{Installing the packages failed -- see
above; Check environment shows what is missing.}), wird der Interpreter der neuen Umgebung
trotzdem eingetragen. Häufige Ursachen sind eine fehlende Internetverbindung, ein
Proxy oder eine Python-Version, für die es noch keine fertigen Pakete gibt. Beheben Sie die
Ursache und installieren Sie die fehlenden Pakete wie im nächsten Abschnitt beschrieben nach.

\subsection{Umgebung prüfen mit \enquote{Check environment}}
\label{sec:check-environment}\index{Check environment}

\gui{Check environment} startet den Interpreter, der \emph{gerade im Feld} steht (auch wenn
er noch nicht mit \gui{OK} übernommen ist), und sucht die Pakete, die die Plugins brauchen.
Der Bericht hat diese Form:

\begin{lstlisting}
Python 3.13.1  ok
v numpy  2.3.4
v raytatouille>=0.9  0.9.2
x pathtracer>=0.3  ModuleNotFoundError: No module named 'pathtracer'

Install the missing packages with:
"/Users/anna/python-env/bin/python" -m pip install "pathtracer>=0.3"
\end{lstlisting}

(Im Programm stehen statt \code{v} und \code{x} ein Häkchen bzw.\ ein Kreuz; Versionen und
Pfade sind Beispiele.) Die erste Zeile nennt die Python-Version; statt \enquote{ok} steht
dort ein Fehler, wenn der Interpreter nicht startet (\enquote{\dots{} cannot be started}) oder
zu alt ist (\enquote{Python 3.9.6 is too old (3.10 or newer needed)}). Für jedes Paket folgt
eine Zeile mit der Anforderung und der gefundenen Version bzw.\ dem Fehler -- auch eine zu
alte Version wird gemeldet (\enquote{version 0.8.1 installed, >=0.9 needed}). Fehlt etwas,
steht am Ende der vollständige \code{pip}-Befehl, den Sie in ein Terminal (Windows:
Eingabeaufforderung oder PowerShell) kopieren können.

\subsection{Eine eigene Umgebung verwenden}

Sie können auch eine bereits vorhandene Umgebung benutzen, etwa eine \code{conda}-Umgebung
oder ein venv, in dem Sie Ihre eigenen Skripte entwickeln:

\begin{schritte}
  \item Installieren Sie die Pakete dort selbst, z.\,B.:
\begin{lstlisting}
python -m pip install numpy scipy "raytatouille>=0.9" "rtt-fusion>=0.2" "photonpower-fieldtracing>=0.2" "fouriergrating>=0.18" "pathtracer>=0.3"
\end{lstlisting}
  \item Tragen Sie unter \menu{Plugins,Settings...} im Feld \gui{Python interpreter} den
    vollständigen Pfad des Interpreters dieser Umgebung ein.
  \item Prüfen Sie mit \gui{Check environment} und bestätigen Sie mit \gui{OK}.
\end{schritte}

Pakete, die nur einzelne Plugins brauchen, können Sie weglassen, wenn Sie diese Plugins nicht
verwenden; die übrigen Plugins laufen trotzdem.

\begin{tipp}
Ein Rechenkern lässt sich jederzeit getrennt aktualisieren:
\code{"<Interpreter>" -m pip install --upgrade raytatouille}. Danach \menu{Plugins,Reload
plugins} wählen, damit die Plugins mit der neuen Version neu starten.
\end{tipp}

Die übrigen Elemente des Dialogs \gui{Plugin settings} -- installierte Plugins, weitere
Plugin-Ordner, vertrauenswürdige Herausgeber, Entwicklermodus und die Liste
\gui{Problems} -- beschreibt \autoref{kap:plugin-verwaltung}.

\section{Der erste Start}\label{sec:erster-start}
\index{erster Start}

Beim ersten Start öffnet sich das Hauptfenster in einer Größe von $1440 \times 900$ Pixeln
mit einem leeren Dokument \enquote{Untitled}:

\begin{itemize}
  \item links die Werkzeugspalte (\gui{Tools}) mit etwa einem Zehntel der Fensterbreite,
  \item in der Mitte die vier Ansichten Front, Top, Left und Perspective,
  \item rechts oben das Dock \gui{Properties} mit \enquote{Nothing selected}, darunter
    die Docks \gui{Layers}, \gui{Model}, \gui{Results} und \gui{Plugin} als Reiter,
  \item unten die Befehlszeile mit dem Text \enquote{Command:} und die Statusleiste.
\end{itemize}

Gleichzeitig lädt das Programm die mitgelieferten Plugins und startet jedes mit dem
eingestellten Python-Interpreter. Solange ein Plugin startet, steht in seinem Untermenü im
Menü \menu{Plugins} \enquote{(starting \dots)}. Ist die Python-Umgebung noch nicht
eingerichtet, scheitert der Start; das Untermenü zeigt dann \enquote{(not running -- see the
plugin log)}, und die Fehlermeldung steht im Dock \gui{Plugin log} (\menu{View,Plugin log})
und in der Statusleiste. Das ist beim allerersten Start normal: Richten Sie die Umgebung wie
in \autoref{sec:create-environment} beschrieben ein.

\paragraph{Ein kurzer Funktionstest.}
\begin{schritte}
  \item Wählen Sie \menu{Help,About Raytatouille Fusion} (auf dem Mac steht der Eintrag im
    Programmmenü \gui{Raytatouille Fusion} ganz links in der Menüleiste). Das Fenster nennt die Programmversion, die Version des
    NURBS-Kerns Nurbforge mit dem Ergebnis seines Selbsttests (\enquote{self-test passed})
    und die Qt-Version.
  \item Wählen Sie \menu{Plugins,Settings...} $\to$ \gui{Check environment}. Alle Zeilen
    sollten ein Häkchen haben.
  \item Öffnen Sie ein Beispiel (\autoref{sec:beispiele}), etwa \datei{demo-parts.step} --
    zwölf Volumenkörper erscheinen in allen vier Ansichten.
\end{schritte}

\begin{hinweis}
Hat das Programm beim vorigen Mal einen Absturzbericht geschrieben, erscheint direkt nach dem
Start ein Hinweis darauf (\autoref{sec:absturzberichte}).
\end{hinweis}

\section{Updates}\label{sec:updates}
\index{Updates}

Das Programm sucht \emph{nicht} von selbst nach Updates. Wählen Sie dazu im Menü
\menu{Help} den Eintrag \gui{Check for updates...}. Das Programm fragt dann

\begin{itemize}
  \item bei GitHub nach der neuesten veröffentlichten Version von \produkt{} und
  \item für jedes Plugin, das in seiner \datei{plugin.json} eine Update-Adresse angibt, dort
    nach dessen neuester Version.
\end{itemize}

Während der Abfrage steht \enquote{Checking for updates \dots} in der Statusleiste. Danach
zeigt ein Fenster \gui{Updates} je Eintrag eine Zeile:

\begin{itemize}
  \item \enquote{Raytatouille Fusion 1.0.0 is up to date} -- nichts zu tun;
  \item \enquote{\textbf{Raytatouille Fusion}: 1.1.0 available (installed 1.0.0) --
    download} -- ein Klick auf \gui{download} öffnet die Seite der neuen Version im Browser;
  \item \enquote{\dots: check failed (\dots)} -- die Abfrage ist gescheitert, etwa ohne
    Internetverbindung.
\end{itemize}

Installiert wird das Update von Hand, wie in den Abschnitten zu den einzelnen Systemen
beschrieben. Bei der Abfrage übermittelt das Programm nur seine Versionsnummer (in der
Kennung \enquote{RaytatouilleFusion/<Version>}).

Die Rechenkerne in der Python-Umgebung werden davon nicht erfasst; sie aktualisieren Sie mit
\code{pip} (siehe Tipp in \autoref{sec:python-umgebung}).

\section{Absturzberichte}\label{sec:absturzberichte}
\index{Absturzbericht}

Stürzt das Programm ab, schreibt es einen Absturzbericht \emph{auf Ihren Rechner}. Es wird
nichts verschickt. Der Bericht ist eine Textdatei mit
einem Namen der Form\\ \datei{crash-<Datum>-<Uhrzeit>-<Prozessnummer>.txt} im Ordner \datei{crash-reports} des
Datenordners (\autoref{sec:einstellungen}) und enthält:

\begin{itemize}
  \item Programmversion, Startzeit (UTC), Betriebssystem und Prozessorarchitektur,
  \item die Ursache (z.\,B.\ \enquote{SIGSEGV (invalid memory access)} oder
    \enquote{std::terminate (uncaught exception)}),
  \item den Aufrufstapel (\enquote{Backtrace}) zum Zeitpunkt des Absturzes.
\end{itemize}

Beim nächsten Start meldet ein Fenster: \enquote{Raytatouille Fusion closed unexpectedly last
time. A crash report was saved on this computer (it is not sent anywhere)} mit dem Pfad des
Berichts. \gui{Show report} öffnet die Datei, \gui{OK} schließt den Hinweis. Jeder Bericht wird
nur einmal gemeldet.

Wenn Sie den Fehler melden möchten, legen Sie den Bericht einer Meldung auf
\url{https://github.com/PhotonPower/Raytatouille-Fusion/issues} bei und beschreiben Sie, was Sie
zuletzt getan haben. Der Bericht enthält keine Modelldaten.

\begin{hinweis}
Abstürze eines Plugins sind keine Abstürze des Programms: Das Plugin läuft in einem eigenen
Python-Prozess. Seine Fehlermeldung steht im Dock \gui{Plugin log}, das sich bei einer
fehlgeschlagenen Rechnung von selbst öffnet.
\end{hinweis}

\section{Wo Einstellungen und Daten gespeichert werden}\label{sec:einstellungen}
\index{Einstellungen}\index{Datenordner}

Das Programm merkt sich seine Einstellungen über das Beenden hinaus und auch über Updates
hinweg. Gespeichert werden:

\begin{referenz}{>{\raggedright\arraybackslash}p{5.2cm}X}
\toprule
\textbf{Einstellung} & \textbf{Inhalt} \\
\midrule
\endhead
Fenstergröße und Anordnung der Docks & beim Beenden gespeichert. Kommt mit einer neuen
  Version ein neues Dock hinzu, wird die alte Anordnung einmal verworfen, damit kein Dock
  verborgen bleibt. \\
Werkzeugspalte & welche Gruppen auf- oder zugeklappt sind \\
Fang & welche Fangarten eingeschaltet sind und ob der Fang insgesamt an ist
  (\gui{Snap}, \keys{F3}) \\
Plugins & Python-Interpreter, weitere Plugin-Ordner, Entwicklermodus, die unter
  \menu{Plugins,Show in the side panels} ausgeblendeten Plugins \\
\bottomrule
\end{referenz}

\emph{Nicht} in den Einstellungen, sondern im Dokument stehen die Parameter der Plugins, der
Workflow, die Darstellung der Ebenen usw.\ (\autoref{kap:dateien}).

\paragraph{Ablageorte.} Die Einstellungen liegen an der für das System üblichen Stelle; der
\emph{Datenordner} enthält Absturzberichte, installierte Plugins und die vorgeschlagene
Python-Umgebung.

\begin{description}
  \item[macOS] Einstellungen:
\begin{lstlisting}
~/Library/Preferences/io.github.photonpower.Raytatouille Fusion.plist
\end{lstlisting}
    Datenordner:
\begin{lstlisting}
~/Library/Application Support/PhotonPower/Raytatouille Fusion
\end{lstlisting}
  \item[Windows] Einstellungen in der Registrierung unter
\begin{lstlisting}
HKEY_CURRENT_USER\Software\PhotonPower\Raytatouille Fusion
\end{lstlisting}
    Datenordner:
\begin{lstlisting}
%APPDATA%\PhotonPower\Raytatouille Fusion
\end{lstlisting}
  \item[Linux] Einstellungen:
\begin{lstlisting}
~/.config/PhotonPower/Raytatouille Fusion.conf
\end{lstlisting}
    Datenordner:
\begin{lstlisting}
~/.local/share/PhotonPower/Raytatouille Fusion
\end{lstlisting}
\end{description}

Im Datenordner liegen:
\begin{itemize}
  \item \datei{crash-reports} -- die Absturzberichte,
  \item \datei{plugins} -- mit \menu{Plugins,Install plugin...} installierte Plugins
    (\autoref{kap:plugin-verwaltung}),
  \item \datei{python-env} -- die mit \gui{Create environment...} angelegte Umgebung, sofern
    Sie den vorgeschlagenen Ordner übernommen haben.
\end{itemize}

\paragraph{Übernahme aus Nurbforge CAD.} Beim ersten Start übernimmt \produkt{} einmalig die
Einstellungen des Vorgängers \enquote{Nurbforge CAD} (Fensteranordnung, Werkzeugspalte, Fang,
Plugin-Einstellungen). Bereits vorhandene eigene Einstellungen werden dabei nicht
überschrieben.

\paragraph{Zurücksetzen.} Um mit den Werkseinstellungen zu beginnen -- etwa wenn ein Dock
verschwunden ist und sich über das Menü \menu{View} nicht zurückholen lässt --, beenden Sie
das Programm und löschen die Einstellungsdatei (unter Windows den Registrierungsschlüssel).
Der Datenordner bleibt davon unberührt.

\section{Deinstallation}\label{sec:deinstallation}
\index{Deinstallation}

\begin{description}
  \item[macOS:] Ziehen Sie \datei{RaytatouilleFusion} aus dem Ordner Programme in den
    Papierkorb.
  \item[Windows:] \gui{Einstellungen} $\to$ \gui{Apps} $\to$ \gui{Installierte Apps} $\to$
    Eintrag \gui{Raytatouille Fusion} $\to$ \gui{Deinstallieren}. Dabei werden auch die
    Verknüpfungen der Dateiendungen \datei{.rttf} und \datei{.nfcad} entfernt.
  \item[Linux:] \code{sudo dnf remove raytatouille-fusion}.
\end{description}

Ihre Einstellungen, der Datenordner (mit Absturzberichten, installierten Plugins und einer
dort angelegten Python-Umgebung) und natürlich Ihre Dokumente bleiben erhalten. Wollen Sie
alles entfernen, löschen Sie zusätzlich die in \autoref{sec:einstellungen} genannten Orte
sowie eine an anderer Stelle angelegte Python-Umgebung (einfach deren Ordner).

\section{Umgebungsvariablen für Sonderfälle}\label{sec:umgebungsvariablen}
\index{Umgebungsvariablen}

Für den Normalbetrieb sind keine Umgebungsvariablen nötig. In Sonderfällen helfen diese:

\begin{referenz}{>{\raggedright\arraybackslash}p{4.4cm}X}
\toprule
\textbf{Variable} & \textbf{Wirkung} \\
\midrule
\endhead
\code{NFCAD\_RHI\_BACKEND} & Grafikschnittstelle der Ansichten: \code{opengl}, \code{vulkan},
  \code{metal}, \code{d3d11} oder \code{d3d12}. Voreingestellt sind Metal (macOS),
  Direct3D~11 (Windows) und OpenGL (Linux). Nützlich bei Treiberproblemen, z.\,B.\
  \code{opengl} in virtuellen Maschinen oder \code{vulkan} unter Linux. \\
\code{NFCAD\_PYTHON} & Python-Interpreter der Plugins, falls unter \menu{Plugins,Settings...}
  keiner eingetragen ist. \\
\code{NFCAD\_PLUGIN\_PATH} & zusätzliche Plugin-Ordner (mehrere getrennt durch \code{:},
  unter Windows durch \code{;}). \\
\code{NFCAD\_DEVELOPER\_MODE} & mit dem Wert \code{1}: Entwicklermodus, lädt auch nicht
  signierte Plugins (\autoref{kap:plugin-verwaltung}). \\
\code{QT\_QPA\_PLATFORM} & unter Linux \code{xcb} erzwingt X11 statt Wayland. \\
\bottomrule
\end{referenz}

Beispiel unter Windows (PowerShell):
\begin{lstlisting}
$env:NFCAD_RHI_BACKEND = "d3d12"
& "C:\Program Files\Raytatouille Fusion\bin\RaytatouilleFusion.exe"
\end{lstlisting}
